\PassOptionsToPackage{table,xcdraw}{xcolor}
\documentclass[aspectratio=169]{beamer}

\usetheme{Madrid}
\usecolortheme{default}
\definecolor{ucbblue}{RGB}{0, 51, 102}

\setbeamercolor{structure}{fg=ucbblue}
\setbeamercolor*{palette primary}{fg=white,bg=ucbblue}
\setbeamercolor*{palette secondary}{fg=white,bg=ucbblue!80}
\setbeamercolor*{palette tertiary}{fg=white,bg=ucbblue!90}
\setbeamercolor{title}{fg=white, bg=ucbblue}
\setbeamercolor{frametitle}{fg=white, bg=ucbblue}
\setbeamercolor{item}{fg=ucbblue}

\usepackage[T1]{fontenc}
\usepackage[utf8]{inputenc}
\usepackage{tgpagella}
\usefonttheme{serif}
\usepackage[spanish, es-noshorthands]{babel}
\usepackage{amsmath, amssymb, bm}
\usepackage{graphicx}
\usepackage[most]{tcolorbox}
\usepackage{booktabs}
\usepackage{tabularx}
\usepackage{tikz}
\usepackage[american]{circuitikz}
\usetikzlibrary{shapes.geometric, arrows.meta, positioning, calc}

\title[Par Diferencial DC]{Transferencia de Pequeña Señal y Análisis DC del Par Diferencial BJT}
\subtitle{IMT-221 Circuitos Electrónicos III — Semana 7, Sesión 2}
\author[B. Quiroga Turdera]{M.Sc. Ing. Bernardo Quiroga Turdera}
\institute[UCB Tarija]{Universidad Católica Boliviana ``San Pablo'' — Sede Tarija}
\date{Semana 7}

\begin{document}

\begin{frame}
    \titlepage
\end{frame}

\begin{frame}{Objetivos de la Sesión}
    \begin{tcolorbox}[colback=ucbblue!5, colframe=ucbblue, title=Transición Conceptual]
        El salto conceptual principal de hoy es pasar de \textbf{una sola rama BJT} (aislada) a \textbf{dos ramas BJT acopladas por una restricción de corriente común}[cite: 16].
    \end{tcolorbox}
    
    \vspace{0.4cm}
    \textbf{Hoja de Ruta:}
    \begin{enumerate}
        \item Consolidación del Método Incremental (CE, CC, CB)[cite: 16].
        \item Estructura Física del Par Diferencial[cite: 16].
        \item Derivación Matemática Paso a Paso del Estado DC Balanceado[cite: 16].
        \item Direccionamiento de Corriente (Current Steering)[cite: 16].
        \item Puente Analítico hacia el Hito 2[cite: 16].
    \end{enumerate}
\end{frame}

\begin{frame}{Revisión: Un Método, Tres Topologías}
    No existen tres sets de fórmulas desconectadas. Existe \textbf{un solo método incremental} invariable[cite: 16]:
    
    \vspace{0.3cm}
    \begin{center}
    \resizebox{0.95\textwidth}{!}{
    \begin{tikzpicture}[auto, thick]
        \node[draw, fill=blue!10, rounded corners, align=center, text width=1.5cm] (q) {1. Punto Q};
        \node[draw, fill=blue!10, rounded corners, align=center, text width=2cm, right=0.3cm of q] (param) {2. Parámetros \\ ($g_m, r_\pi, r_e$)};
        \node[draw, fill=blue!10, rounded corners, align=center, text width=2cm, right=0.3cm of param] (ac) {3. Equivalente AC};
        \node[draw, fill=blue!10, rounded corners, align=center, text width=2.4cm, right=0.3cm of ac] (eq) {4. Leyes \\ (KVL/KCL)};
        \node[draw, fill=green!10, rounded corners, align=center, text width=1.8cm, right=0.3cm of eq] (res) {$A_v, R_{in}, R_{out}$};
        
        \draw[->] (q) -- (param); \draw[->] (param) -- (ac); \draw[->] (ac) -- (eq); \draw[->] (eq) -- (res);
    \end{tikzpicture}
    }
    \end{center}
    
    \vspace{0.4cm}
    \textbf{Lo que cambia es puramente topológico:}
    \begin{itemize}
        \item \textbf{CE:} Entrada Base, Salida Colector, \textbf{Emisor Común (Tierra AC)}. Ecuación: $v_o = -g_m v_i R_C$.
        \item \textbf{CC:} Entrada Base, Salida Emisor, \textbf{Colector Común (Tierra AC)}. Ecuación: $v_o = i_e R_E$.
        \item \textbf{CB:} Entrada Emisor, Salida Colector, \textbf{Base Común (Tierra AC)}. Ecuación: $v_o = -(-g_m v_i) R_C$.
    \end{itemize}
\end{frame}

\begin{frame}{Estructura del Par Diferencial Canónico}
    \begin{columns}
        \begin{column}{0.45\textwidth}
            \centering
            \begin{circuitikz}[scale=0.75, transform shape, thick]
                \draw (0,0) node[npn](Q1){}; \node[left=0.2cm] at (Q1.B) {$v_1$};
                \draw (3,0) node[npn, xscale=-1](Q2){}; \node[right=0.2cm] at (Q2.B) {$v_2$};
                \draw (Q1.C) to[R, l={$R_C$}, i={$I_{C1}$}] (0,2.5) -- (3,2.5) to[R, l_={$R_C$}, i_={$I_{C2}$}] (Q2.C);
                \draw (1.5,2.5) -- (1.5,3) node[above]{$+V_{CC}$};
                \draw (Q1.E) -- (0,-1.5) -- (3,-1.5) -- (Q2.E);
                \draw (1.5,-1.5) node[circle, fill, inner sep=1.5pt]{} node[above right]{Nodo $E$} to[I, l={$I_T$}] (1.5,-3.0) node[below]{$-V_{EE}$};
                \draw (Q1.B) to[short, -o] (-1,0);
                \draw (Q2.B) to[short, -o] (4,0);
                \draw (Q1.C) to[short, *-o] (-1,1) node[left]{$v_{C1}$};
                \draw (Q2.C) to[short, *-o] (4,1) node[right]{$v_{C2}$};
            \end{circuitikz}
        \end{column}
        \begin{column}{0.52\textwidth}
            \textbf{Anatomía de la etapa acoplada[cite: 16]:}
            \begin{itemize}
                \item \textbf{Dos ramas} ($Q_1$ y $Q_2$) idealmente idénticas (\textit{matched}).
                \item Dos nodos de entrada independientes ($v_1, v_2$).
                \item Resistores de colector simétricos ($R_C$).
                \item \textbf{Restricción de Acoplamiento:} Los emisores están unidos y alimentados por una sola fuente de cola ($I_T$)[cite: 16].
            \end{itemize}
            
            \vspace{0.2cm}
            \begin{alertblock}{La Ecuación Maestra DC}
                Por Ley de Corrientes de Kirchhoff en el Nodo E:
                \begin{equation*}
                    \mathbf{I_{E1} + I_{E2} = I_T} \quad \text{[cite: 16]}
                \end{equation*}
            \end{alertblock}
        \end{column}
    \end{columns}
\end{frame}

\begin{frame}{Etapa 2: Espejo de Corriente (Fuente de Cola)}
    \begin{columns}[c]
        \begin{column}{0.4\textwidth}
            \centering
            \begin{circuitikz}[scale=0.9, transform shape, thick]
                % Q4 (Salida) a la izquierda, base apuntando a la derecha
                \draw (0,0) node[npn, xscale=-1](Q4){}; 
                \node at (-0.6, -0.2) {Q4};
                
                % Q3 (Referencia) a la derecha, base apuntando a la izquierda
                \draw (2.5,0) node[npn](Q3){}; 
                \node at (3.1, -0.2) {Q3};
                
                % Conexión de emisores a -9V
                \draw (Q4.E) -- (0,-1.2) node[below]{$-9\,\text{V}$};
                \draw (Q3.E) -- (2.5,-1.2) node[below]{$-9\,\text{V}$};
                
                % Conexión de bases (ahora enfrentadas) y nodo V_B
                \draw (Q4.B) -- (1.25,0) node[circle,fill,inner sep=1.5pt](Bmid){} -- (Q3.B);
                \node[below] at (Bmid) {$V_B$};
                
                % Conexión diódica de Q3 (Colector a Base) con ángulos rectos
                \draw (Q3.C) node[circle,fill,inner sep=1.5pt]{} -| (Bmid);
                
                % Rama de referencia (R_REF hacia 0V)
                \draw (2.5,2.5) node[above]{$0\,\text{V}$} to[R, l={$R_{REF}$}, a={$8.2\,\text{k}\Omega$}] (Q3.C);
                
                % Rama de salida (I_T hacia Nodo E)
                \draw (Q4.C) -- (0,2) node[above]{Hacia Nodo E};
                \draw[<-, red, thick] (0,1.3) -- (0,0.8) node[midway, left]{$I_T$};
            \end{circuitikz}
        \end{column}
        \begin{column}{0.55\textwidth}
            \textbf{Predicción Analítica[cite: 18]:}
            En Q3 (Diodo conectado), asumimos $V_{BE} \approx 0.7\,\text{V}$.
            El voltaje en el nodo Base/Colector es:
            \begin{equation*}
                V_B \approx -9\,\text{V} + 0.7\,\text{V} = \mathbf{-8.3\,\text{V}} \quad \text{[cite: 18]}
            \end{equation*}
            Aplicando Ley de Ohm en $R_{REF}$[cite: 18]:
            \begin{equation*}
                I_{REF} = \frac{0\,\text{V} - (-8.3\,\text{V})}{8.2\,\text{k}\Omega} \approx \mathbf{1.01\,\text{mA}} \quad \text{[cite: 18]}
            \end{equation*}
            Para $\beta$ finito, $I_{REF} = I_O (1 + 2/\beta)$[cite: 18]. 
            $$ \mathbf{I_T \approx \frac{1.01\,\text{mA}}{1 + 2/\beta}} \quad \text{[cite: 18]} $$
        \end{column}
    \end{columns}
\end{frame}

\begin{frame}{Análisis DC Paso a Paso: Corrientes de Emisor}
    Derivación explícita cuando las entradas son iguales ($v_1 = v_2 = V_{CM}$) y los transistores son idénticos[cite: 16]:
    
    \vspace{0.3cm}
    \textbf{1. Simetría de Corrientes de Emisor:}
    La tensión base-emisor en ambas ramas es:
    \begin{align*}
        v_{BE1} &= v_1 - v_E \\
        v_{BE2} &= v_2 - v_E
    \end{align*}
    Como $v_1 = v_2$, entonces $v_{BE1} = v_{BE2}$. En transistores idénticos, mismos $v_{BE}$ generan mismas corrientes: $I_{E1} = I_{E2}$[cite: 16].
    
    \vspace{0.2cm}
    Aplicando KCL en el nodo común:
    \begin{equation*}
        I_{E1} + I_{E2} = I_T \implies 2 I_{E1} = I_T \implies \mathbf{I_{E1} = I_{E2} = \frac{I_T}{2}} \quad \text{[cite: 16]}
    \end{equation*}
\end{frame}

\begin{frame}{Análisis DC Paso a Paso: Colectores}
    \textbf{2. Corrientes de Colector (Punto Q)[cite: 16]:}
    Sabiendo que $I_C = \alpha I_E$, donde $\alpha = \frac{\beta}{\beta+1}$[cite: 16]:
    \begin{equation*}
        \mathbf{I_{C1} = I_{C2} = \alpha \frac{I_T}{2} \approx \frac{I_T}{2}} \quad \text{(Aproximación para $\beta$ grande)} \quad \text{[cite: 16]}
    \end{equation*}
    
    \vspace{0.3cm}
    \textbf{3. Tensiones de Colector (Salida DC)[cite: 16]:}
    Aplicamos KVL desde la fuente $V_{CC}$ descendiendo por el resistor $R_C$:
    \begin{align*}
        v_{C1} &= V_{CC} - I_{C1}R_C \\
        v_{C2} &= V_{CC} - I_{C2}R_C
    \end{align*}
    Como $I_{C1} = I_{C2}$, entonces $\mathbf{v_{C1} = v_{C2}}$[cite: 16]. La salida diferencial quiescente es $V_{OD,Q} = v_{C2} - v_{C1} = 0\,\text{V}$[cite: 16].
\end{frame}

\begin{frame}{Análisis DC Paso a Paso: Emisor y Verificación}
    \textbf{4. Voltaje de Emisor y $V_{CE}$[cite: 16]:}
    Asumiendo $V_{BE} \approx 0.7\,\text{V}$ en conducción:
    \begin{equation*}
        v_E = v_1 - V_{BE} = V_{CM} - 0.7\,\text{V} \quad \text{[cite: 16]}
    \end{equation*}
    El voltaje Colector-Emisor en el punto $Q$ para ambas ramas es:
    \begin{equation*}
        \mathbf{V_{CE} = v_C - v_E} \quad \text{[cite: 16]}
    \end{equation*}
    
    \vspace{0.3cm}
    \textbf{5. Verificación de Región Activa Directa[cite: 16]:}
    Debe verificarse rigurosamente que las uniones Base-Colector estén polarizadas en inversa:
    \begin{equation*}
        V_{BC} = v_B - v_C < 0 \quad \text{[cite: 16]}
    \end{equation*}
    Solo así se garantiza que los transistores no han entrado en saturación[cite: 16].
\end{frame}

\begin{frame}{Ejemplo Numérico: Corrientes de Equilibrio}
    \textbf{Datos (Ejemplo Pedagógico):} $V_{CC}=5\,\text{V}$, $I_T=1.0\,\text{mA}$, $R_C=3.3\,\text{k}\Omega$, $v_1=v_2=0\,\text{V}$, $\beta=100$, $V_{BE}\approx0.7\,\text{V}$[cite: 16].
    
    \vspace{0.3cm}
    \textbf{Paso 1 - Corrientes de Emisor:}
    \begin{equation*}
        I_{E1} = I_{E2} = \frac{I_T}{2} = \frac{1.0\,\text{mA}}{2} = \mathbf{0.500\,\text{mA}} \quad \text{[cite: 16]}
    \end{equation*}
    
    \vspace{0.3cm}
    \textbf{Paso 2 - Factor $\alpha$ y Corrientes de Colector:}
    \begin{equation*}
        \alpha = \frac{\beta}{\beta+1} = \frac{100}{101} \approx \mathbf{0.9901} \quad \text{[cite: 16]}
    \end{equation*}
    \begin{equation*}
        I_{C1} = I_{C2} = \alpha I_E = (0.9901)(0.500\,\text{mA}) \approx \mathbf{0.495\,\text{mA}} \quad \text{[cite: 16]}
    \end{equation*}
\end{frame}

\begin{frame}{Ejemplo Numérico: Tensiones de Colector y Emisor}
    Continuando con el circuito ($I_C \approx 0.495\,\text{mA}$, $V_{CC}=5\,\text{V}$, $R_C=3.3\,\text{k}\Omega$, $v_1=0\,\text{V}$, $V_{BE}\approx0.7\,\text{V}$)[cite: 16]:
    
    \vspace{0.3cm}
    \textbf{Paso 3 - Tensiones de Colector:}
    \begin{equation*}
        v_C = V_{CC} - I_C R_C = 5\,\text{V} - (0.495\,\text{mA})(3.3\,\text{k}\Omega) = 5 - 1.6335 = \mathbf{3.366\,\text{V}} \quad \text{[cite: 16]}
    \end{equation*}
    
    \vspace{0.3cm}
    \textbf{Paso 4 - Tensión de Emisor y $V_{CE}$:}
    \begin{equation*}
        v_E = v_1 - V_{BE} = 0\,\text{V} - 0.7\,\text{V} = \mathbf{-0.700\,\text{V}} \quad \text{[cite: 16]}
    \end{equation*}
    \begin{equation*}
        V_{CE} = v_C - v_E = 3.366\,\text{V} - (-0.700\,\text{V}) = \mathbf{4.066\,\text{V}} \quad \text{[cite: 16]}
    \end{equation*}
\end{frame}

\begin{frame}{Ejemplo Numérico: Verificación de Región y Punto Q}
    \textbf{Paso 5 - Verificación Final (Región de Operación):}
    Evaluamos la unión Base-Colector para asegurar que no existe saturación:
    \begin{equation*}
        V_{BC} = v_B - v_C = 0\,\text{V} - 3.366\,\text{V} = \mathbf{-3.366\,\text{V}} \quad \text{[cite: 16]}
    \end{equation*}
    
    \vspace{0.3cm}
    \begin{tcolorbox}[colback=green!5, colframe=green!50!black, title=\textbf{Conclusión Matemática}]
        Como $V_{BC} < 0$, se confirma que las uniones están polarizadas en inversa y los transistores operan correctamente en \textbf{Activa Directa}[cite: 16]. \\
        El Punto de Operación estático para cada rama es:
        \begin{equation*}
            \mathbf{Q_1 = Q_2 \approx (0.495\,\text{mA},\, 4.066\,\text{V})} \quad \text{[cite: 16]}
        \end{equation*}
    \end{tcolorbox}
\end{frame}

\begin{frame}{Direccionamiento de Corriente (Current Steering)}
    Definimos el voltaje diferencial como: $\mathbf{v_d = v_1 - v_2}$[cite: 16]. ¿Qué pasa cualitativamente si $v_d \neq 0$?[cite: 16]
    
    \vspace{0.1cm}
    \begin{table}
        \centering
        \small
        \renewcommand{\arraystretch}{1.3}
        \begin{tabularx}{0.95\textwidth}{|>{\raggedright\arraybackslash}p{2.5cm}|>{\centering\arraybackslash}X|>{\centering\arraybackslash}X|>{\centering\arraybackslash}X|>{\centering\arraybackslash}X|}
        \hline
        \rowcolor{ucbblue!20}
        \textbf{Condición} & $\mathbf{I_{C1}}$ & $\mathbf{I_{C2}}$ & $\mathbf{v_{C1}}$ & $\mathbf{v_{C2}}$ \\ \hline
        $v_1 = v_2$ ($v_d = 0$) & Iguales & Iguales & Iguales & Iguales \\ \hline
        $\mathbf{v_1 > v_2}$ ($v_d > 0$) & Sube ($\uparrow$) & Baja ($\downarrow$) & Cae ($\downarrow$) & Sube ($\uparrow$) \\ \hline
        $\mathbf{v_1 < v_2}$ ($v_d < 0$) & Baja ($\downarrow$) & Sube ($\uparrow$) & Sube ($\uparrow$) & Cae ($\downarrow$) \\ \hline
        \end{tabularx}
    \end{table}
    
    \vspace{0.1cm}
    \begin{alertblock}{La Clave de la Interacción: KVL y KCL}
        La corriente total $I_T$ impone que lo que gana $Q_1$, lo pierde $Q_2$ ($I_{C1} \uparrow \implies I_{C2} \downarrow$)[cite: 16]. Por KVL ($v_C = V_{CC} - I_C R_C$), el voltaje del colector se mueve en \textbf{dirección opuesta} a la corriente del colector[cite: 16].
    \end{alertblock}
\end{frame}

\begin{frame}{El Diseño Real: Sensibilidad, "Mismatch" y el Hito 2}
    En hardware real, incluso con $v_1 = v_2$, casi nunca observamos $v_{C1} = v_{C2}$ perfectos. ¿Por qué?[cite: 16]
    \begin{itemize}
        \item \textbf{Mismatch de transistores:} $V_{BE}$ y $\beta$ difieren ligeramente por fabricación[cite: 16].
        \item \textbf{Mismatch de Resistores:} Las tolerancias de $R_{C1}$ y $R_{C2}$ causan caídas distintas[cite: 16].
        \item Todo esto genera un \textbf{Voltaje de Offset} en la salida[cite: 16].
    \end{itemize}
    
    \vspace{0.2cm}
    \textbf{Preparación para el Laboratorio Hito 2 (Semana 8, Sesión 2)[cite: 16]:}
    \begin{itemize}
        \item El proyecto requerirá medir corrientes con una resistencia Shunt ($V_{sh} = I_{sense} \cdot R_{sh}$)[cite: 17].
        \item La cadena de señal será: $\mathbf{I_{load} \to R_{sh} \to v_d \to \text{Par Diferencial} \to v_{od}}$[cite: 17].
        \item \textbf{NO armar ni comprar} sin revisar la "Guía Previa" oficial; utilizaremos un \textbf{Espejo de Corriente} en la cola para mejorar el CMRR[cite: 17].
    \end{itemize}
\end{frame}

\end{document}
